\section*{第 \ref{ch:g5:areas} 章——周长与面积}

\begin{solution}{exo:g5:areas:1}
\omterm{def:g4:measure:area}{面积}：$8 \times 5 = 40$ cm$^2$。\omterm{def:g4:measure:perimeter}{周长}：
$2 \times (8 + 5) = 26$ cm。
\end{solution}

\begin{solution}{exo:g5:areas:2}
$9 \times 9 = 81$ cm$^2$; \quad $12 \times 7 = 84$ m$^2$; \quad
$4.5 \times 6 = 27$ cm$^2$.
\end{solution}

\begin{solution}{exo:g5:areas:3}
宽：$63 \div 9 = 7$ cm。\omterm{def:g4:measure:perimeter}{周长}：$2 \times (9 + 7) = 32$ cm。
\end{solution}

\begin{solution}{exo:g5:areas:4}
例如 $6 \times 4$（\omterm{def:g4:measure:perimeter}{周长} $20$）和 $8 \times 3$
（\omterm{def:g4:measure:perimeter}{周长} $22$）：\omterm{def:g4:measure:area}{面积}相同，都是 $24$，\omterm{def:g4:measure:perimeter}{周长}不同。
\end{solution}

\begin{solution}{exo:g5:areas:5}
$3$ m$^2 = 30\,000$ cm$^2$; \quad $50\,000$ cm$^2 = 5$ m$^2$.
\end{solution}

\begin{solution}{exo:g5:areas:6}
横条：$8 \times 2 = 16$ cm$^2$；竖条：$2 \times 5 = 10$ cm$^2$；
合计 $26$ cm$^2$。
\end{solution}

\begin{solution}{exo:g5:areas:7}
所需木料为 $30 \times 20 - 24 \times 14 = 600 - 336 = 264$ cm$^2$。
\end{solution}

\begin{solution}{exo:g5:areas:8}
$8 \times 6$ 长方形的\omterm{def:g3:division:half}{一半}：$48 \div 2 = 24$ cm$^2$。
\end{solution}

\begin{solution}{exo:g5:areas:9}
草籽：\omterm{def:g4:measure:area}{面积} $7 \times 4 = 28$ m$^2$，费用 $28 \times 2 = 56$。
围栏：\omterm{def:g4:measure:perimeter}{周长} $2 \times (7 + 4) = 22$ m，费用
$22 \times 3 = 66$.
\end{solution}

\begin{solution}{exo:g5:areas:10}
两块长 $50$ cm 的瓷砖拼成一米，沿 $12$ m 的长需
$24$ 块；沿 $2$ m 的宽需 $4$ 块。合计：
$24 \times 4 = 96$ 块。
\end{solution}

\begin{solution}{exo:g5:areas:11}
原来：\omterm{def:g4:measure:perimeter}{周长} $2 \times (4 + 3) = 14$ cm，\omterm{def:g4:measure:area}{面积} $12$ cm$^2$。
变化后（边长 $8$ 和 $6$）：\omterm{def:g4:measure:perimeter}{周长} $28$ cm，\omterm{def:g4:measure:area}{面积} $48$ cm$^2$。
\omterm{def:g4:measure:perimeter}{周长}\emph{翻了一倍}；\omterm{def:g4:measure:area}{面积}则乘以
\emph{四}倍（$2 \times 2$）：长度与\omterm{def:g4:measure:area}{面积}的缩放
方式并不相同。
\end{solution}
